% ----- Consignes exo 56 ----- %
\begin{td-exo}[Exemples de \(L\)-réductions de base]\,\\ % 56 
	On rappelle les énoncés des deux problèmes suivants:

	\begin{problem}{\textsc{Maximum} \(\neq 3\)\textsc{-SAT}}
    \Input{Même données que pour \textsc{3-SAT}}
    \Question{Trouver le maximum de clauses satisfaites telle que dans chaque clause satisfaite il existe au moins un littéral mis à vrai et au moins un littéral mis à faux.}
\end{problem}
	\begin{problem}{\textsc{Maximum Multi-Coupe} (\textsc{Maximum Multi-Cut})}
    \Input{Un multigraphe \(G=(V,E)\) et un entier \(k\)}
    \Question{Trouver un ensemble de sommets \(V_1,V_2 \subseteq V\) tel que \(V_1 \cup V_2 = V\) et que le nombre d'arêtes entre \(V_1\) et \(V_2\) est au moins \(k\).}
\end{problem}

	Faire les \(L\)-réductions sur les problèmes suivants:
	\begin{enumerate}
		\item \(2\)-\textsc{Satisfiabilite Maximale} \(\leq_L\) \textsc{Maximum} \(\neq 3\)-\textsc{SAT}

		\item \textsc{Maximum} \(\neq 3\)-\textsc{SAT} \(\leq_L\) \textsc{Maximum Multi-Coupe}.
		On propose la construction suivante:
		\begin{constr}
			Considérons une instance \(\varphi\) de \textsc{Maximum} \(\neq 3\)-\textsc{SAT} sur \(n\) variables et \(m\) clauses et supposons que chacune des \(m\) clauses de \(\varphi\) contient au moins un littéral.
			\begin{itemize}
				\item Pour toute variable \(x\) de \(\varphi\), la graphe \(G = f(\varphi)\) contient deux sommets \(v_x\) et \(v_{\ol{x}}\); ces sommets sont reliés entre eux par \(2n_x\) arêtes parallèles, où \(n_x\) est le nombre de fois que la variable \(x\) apparaît dans \(\varphi\).
				Nous appelons ces arêtes ``\(v\)-arêtes'' pour indiquer qu'elles sont associées aux variables de \(\varphi\).

				\item Pour chaque clause \(C = (y \lor z \lor w)\) de \(\varphi\), le graphe \(G\) contient un triangle sur les sommets \(v_y\), \(v_z\) et \(v_w\); pour chaque clause \(C = (y\lor z)\) de \(\varphi\), le graphe \(G\) contient deux arêtes parallèles entre les sommets \(v_y\) et \(v_z\).
				Nous allons appeler ces arêtes ``\(c\)-arêtes'' pour indiquer qu'elles sont associées aux clauses de \(\varphi\).
				La fonction \(f\) de la réduction est ainsi spécifiée.
				Il est facile de voir que le graphe ainsi construit est un multigraphe, instance de \textsc{Maximum Multi-Coupe}.
			\end{itemize}
		\end{constr}
	\end{enumerate}
\end{td-exo}

% ----- Solutions exo 56 ----- %
\iftoggle{showsolutions}{ 
	\begin{td-sol}[]\ % 56
		\begin{enumerate}
			\item On considère la transformation suivante:
			\begin{equation*}
				\left(l_1^i, l_2^i\right) \mapsto \left(l_1^i, l_2^i, x\right) 
			\end{equation*}
			Alors, toute solution de \(2\)-\textsc{Satisfiabilite Maximale} est aussi une solution de \textsc{Maximum} \(\neq 3\)-\textsc{SAT}, il suffit de mettre \(x\) à faux. De plus, toute solution de \textsc{Maximum} \(\neq 3\)-\textsc{SAT} est aussi une solution de \(2\)-\textsc{Satisfiabilite Maximale}, il suffit d'ignorer les clauses contenant \(x\). Par conséquent, on a bien une \(L\)-réduction avec \(\alpha_1 = \alpha_2 = 1\).
		\end{enumerate}

		On se propose dans la suite de montrer les relations suivantes:
		\begin{equation*}
			\min \textsc{VC}(3) \prec_L \min \textsc{Dom}(6) \prec_L \min \textsc{Dom}(3) 
		\end{equation*}

		Montrons d'abord que \(\min \textsc{VC}(3) \prec_L \min \textsc{Dom}(6)\).
		Soit \(G\) un graphe.
		On propose la construction suivante:
		\begin{equation*}
			\forall (u, v) \in E(G),
		\end{equation*}
		On construit le triangle \((u, v, w)\).

		Montrons maintenant qu'on a \(\forall k, \textsc{VC}(k) \approx \textsc{Dom}(k)\).

		Si on a une solution à \textsc{VC}\((k)\), le même ensemble de sommets est bien une solution à \textsc{Dom}\((k)\) puisque tout sommet de \(G\) est soit dans la solution, soit adjacent à un sommet de la solution.

		Si on a une solution à \textsc{Dom}\((k)\), on effectue les opérations suivantes:
		\begin{itemize}
			\item Si \(u\) ou \(v\) est dans la solution, on ajoute \(u\) et \(v\) à la solution.
			\item Sinon, \(w\) est dans la solution, on ajoute alors \(u\) ou \(v\) à la solution.
		\end{itemize}

		On a donc égalité des solutions optimales pour les deux problèmes, et on a aussi égalité des solutions approximatives, ce qui montre que \(\textsc{VC}(k) \approx \textsc{Dom}(k)\) avec \(\alpha = \beta = 1\).

		Si on regarde maintenant le degré des graphes, dans un graphe subcubique, un sommet transformé en triangle a un degré au plus égal à 6, d'où la réduction de \(\textsc{VC}(3)\) à \(\textsc{Dom}(6)\).

		Montrons maintenant que \(\min \textsc{Dom}(6) \prec_L \min \textsc{Dom}(3)\).
	\end{td-sol}
}{}
